\scanpage{22}
\begin{center}--- 50 ---\end{center}

Un cas particulier intéressant est celui où la variété abélienne est \textit{ordinaire}, autrement dit a le maximum possible de points d'ordre $p$. Il existe alors un \textit{relèvement canonique} que l'on peut caractériser de diverses manières; dans le cas des courbes elliptiques, ce relèvement avait été signalé en 1933 par H. \textsc{Hasse} (qui se plaçait à un point de vue quelque peu différent).

\begin{center}\textsc{Publications}\end{center}

H. \textsc{Bass}, J. \textsc{Milnor} et J.-P. \textsc{Serre}, \textit{Solution of the congruence subgroup problem for $\mathrm{SL}_n$ $(n\geq3)$ and $\mathrm{Sp}_{2n}$ $(n\geq2)$} (\textit{Publ. Math. I.H.E.S.}, n\textsuperscript{o} 33, 1967, p. 59-137).

J.-P. \textsc{Serre}, \textit{Local Class Field Theory}; \textit{Complex Multiplication} (chap. VI et XIII de J. \textsc{Cassels} et A. \textsc{Fröhlich}, \textit{Algebraic Number Theory}, Acad. Press, New York, 1967).

--- \textit{Sur les groupes de Galois attachés aux groupes $p$-divisibles} (\textit{Proc. Conf. Local Fields}, Driebergen, 1966, Springer-Verlag, p. 118-131).

--- \textit{Commutativité des groupes formels de dimension 1} (\textit{Bull. Sci. math.}, 91, 1967, p. 113-115).

--- \textit{Représentations linéaires des groupes finis} (Paris, Hermann, 1968).

--- \textit{Abelian $l$-adic representations and elliptic curves} (Notes rédigées avec la collaboration de W. \textsc{Kuyk} et J. \textsc{Labute}, W. A. Benjamin Publ., New York, 1968).

\begin{center}\textsc{Missions}\end{center}

Cours à McGill University, Montréal (septembre 1967).

Conférences à l'Université de Montréal (septembre 1967), à Haverford College (octobre 1967), à l'Institute for Advanced Study de Princeton (octobre-décembre 1967) et à l'Université de Grenoble (juin 1968).
\margin{21}
